\section*{Capítulo \ref{ch:b1:expression-control} --- Controle da expressão gênica}

\begin{solution}{exo:b1:expression-control:1}
O \emph{lacI} (\omterm{def:b1:genomes:gene}{gene} do \omterm{def:b1:expression-control:operon}{repressor}), o promotor, o operador, o
\emph{lacZ}, o \emph{lacY}, o \emph{lacA} e o sítio da CAP. Sem lactose:
desligado (\omterm{def:b1:expression-control:operon}{repressor} ligado). Com lactose e com glicose: mal ligado
(\omterm{def:b1:expression-control:operon}{repressor} solto, mas sem AMP cíclico e CAP). Com lactose e sem glicose:
plenamente ligado. Sem nenhum dos dois açúcares: desligado.
\end{solution}

\begin{solution}{exo:b1:expression-control:2}
Um produto que age em trans é uma molécula difusível que pode agir sobre
qualquer cópia do alvo dela na \omterm{def:b1:cell-unit-of-life:cell}{célula}: o \omterm{def:b1:expression-control:operon}{repressor}. Um sítio que age em
cis é uma sequência de \omterm{def:b1:nucleic-acids:chain}{DNA} que afeta apenas os \omterm{def:b1:genomes:gene}{genes} fisicamente ligados
a ela: o operador.
\end{solution}

\begin{solution}{exo:b1:expression-control:3}
Latência de cerca de \qty{36}{min} (de \qty{2}{h} a \qty{2.6}{h});
duplicação de \qty{20}{min} em glicose e de \qty{30}{min} em lactose.
\end{solution}

\begin{solution}{exo:b1:expression-control:4}
\omterm{def:b1:expression-control:enhancer}{Intensificador}: uma sequência reguladora, muitas vezes distante, que
eleva a transcrição de um \omterm{def:b1:genomes:gene}{gene} quando fatores se ligam a ela. \omterm{def:b1:expression-control:enhancer}{Fator de
transcrição}: uma \omterm{def:b1:proteins:peptide}{proteína} com um domínio de ligação ao \omterm{def:b1:nucleic-acids:chain}{DNA} e um domínio
de ativação ou de repressão que regula a transcrição. \omterm{def:b1:genomes:gene}{Gene} de
manutenção: um \omterm{def:b1:genomes:gene}{gene} expresso em toda \omterm{def:b1:cell-unit-of-life:cell}{célula} num nível constante.
\end{solution}

\begin{solution}{exo:b1:expression-control:5}
$I^-$: constitutivo. $O^c$: constitutivo. $I^s$: não induzível.
$I^-/I^+$: induzível (o \omterm{def:b1:expression-control:operon}{repressor} $I^+$ age em trans). $I^s/I^+$: não
induzível ($I^s$ é dominante). $I^+\,O^c\,Z^- / I^+\,O^+\,Z^+$:
induzível --- o único $Z$ funcional fica atrás de um operador normal.
\end{solution}

\begin{solution}{exo:b1:expression-control:6}
Lac: o açúcar (indutor) se liga ao \omterm{def:b1:expression-control:operon}{repressor} e o \emph{solta} do
operador --- as \omterm{def:b1:enzymes:enzyme}{enzimas} são fabricadas quando o \omterm{def:b1:enzymes:enzyme}{substrato} delas está
presente. Trp: o \omterm{def:b1:proteins:aminoacid}{aminoácido} (correpressor) se liga ao \omterm{def:b1:expression-control:operon}{repressor} e o
\emph{habilita} a se ligar --- as \omterm{def:b1:enzymes:enzyme}{enzimas} são fabricadas quando o
produto delas está ausente. As vias catabólicas são ligadas pela entrada
delas, e as anabólicas, desligadas pela saída delas.
\end{solution}

\begin{solution}{exo:b1:expression-control:7}
Em glicose: crescimento normal (a glicose não precisa de CAP). Só em
lactose: quase nenhum crescimento --- sem AMP cíclico, a CAP não
consegue ativar o promotor lac, de modo que a \omterm{def:b1:enzymes:enzyme}{enzima} fica em poucos por
cento do nível induzido, mesmo com o \omterm{def:b1:expression-control:operon}{repressor} solto. Nos dois:
crescimento em glicose e depois uma parada; não há segunda fase. A
\omterm{def:b1:enzymes:enzyme}{enzima} fica baixa em todos os casos.
\end{solution}

\begin{solution}{exo:b1:expression-control:8}
A atenuação precisa que o \omterm{def:b1:gene-expression:ribosome}{ribossomo} esteja traduzindo o mensageiro
enquanto a polimerase ainda o transcreve, de modo que a posição do
\omterm{def:b1:gene-expression:ribosome}{ribossomo} molde o \omterm{def:b1:nucleic-acids:chain}{RNA} atrás da polimerase. Nos eucariotos a transcrição
está no núcleo e a tradução no \omterm{def:b1:eukaryotic-cell:organelle}{citosol}, e a mensagem só é traduzida
depois de processada e exportada.
\end{solution}

\begin{solution}{exo:b1:expression-control:9}
Os \omterm{def:b1:expression-control:enhancer}{intensificadores} agem a distância, em qualquer orientação e de
qualquer lado, de modo que não podem funcionar por serem lidos nem por
posicionarem diretamente a polimerase; eles funcionam ligando fatores
que tocam o promotor por meio de uma alça de \omterm{def:b1:nucleic-acids:chain}{DNA}, para a qual a
distância e a orientação pouco importam.
\end{solution}

\begin{solution}{exo:b1:expression-control:10}
$5000\times 4096\times 4 = \num{8.2e7}$ \omterm{def:b1:water-small-molecules:atp}{ATP}: \qty{0.8}{\%} do orçamento.
Velocidade de crescimento \qty{0.8}{\%} menor: por geração a parcela do
mutante cai pelo fator $2^{-0.008} = 0.9945$. De $1:1$ a $1:99$:
$0.9945^n = 1/99$, $n = \ln 99/0.00555 = 830$ gerações --- algumas
semanas de cultura contínua.
\end{solution}

\begin{solution}{exo:b1:expression-control:11}
Cada tipo celular carrega um conjunto específico de \omterm{def:b1:expression-control:enhancer}{fatores de
transcrição}, que se ligam aos \omterm{def:b1:expression-control:enhancer}{intensificadores} dos \omterm{def:b1:genomes:gene}{genes} dele e recrutam
acetilases que abrem a cromatina deles, enquanto os \omterm{def:b1:genomes:gene}{genes} dos outros
tipos ficam metilados e empacotados em heterocromatina que fator nenhum
alcança. Os fatores mantêm a expressão uns dos outros (uma rede com
estados estáveis), e as marcas de cromatina são copiadas na replicação,
de modo que as filhas herdam tanto os fatores quanto as regiões abertas
e fechadas. O \omterm{def:b1:genomes:genome}{genoma} é o mesmo; o \omterm{def:b1:genomes:genome}{genoma} acessível não é, e a
inacessibilidade é herdada.
\end{solution}

\begin{solution}{exo:b1:expression-control:12}
Os controles de uma bactéria são sintonizados no meio: um açúcar, um
\omterm{def:b1:proteins:aminoacid}{aminoácido}, uma temperatura liga um \omterm{def:b1:expression-control:operon}{repressor} ou um \omterm{prop:b1:expression-control:trp}{fator sigma}, e a
resposta se completa em minutos e se desfaz assim que o sinal
desaparece --- a expressão acompanha o ambiente. A \omterm{def:b1:cell-unit-of-life:cell}{célula} de um corpo
também responde a sinais, mas as decisões principais dela foram tomadas
durante o desenvolvimento e travadas na cromatina: o que ela pode
expressar foi fixado pelos fatores que ela herdou e pelas marcas no \omterm{def:b1:nucleic-acids:chain}{DNA}
dela, alteradas só pela divisão e quase sempre irreversíveis. A bactéria
é uma leitora do presente; a \omterm{def:b1:cell-unit-of-life:cell}{célula} diferenciada tem uma memória que é a
identidade dela.
\end{solution}

\begin{solution}{pb:b1:expression-control:1}
\textbf{1.} $0.5\times 10^{-3}\times 180 = \qty{0.09}{g/L} =
\qty{9e-5}{g/mL}$.
\textbf{2.} $9\times 10^{-5}/1.5\times 10^{-12} = \num{6e7}$ divisões.
\textbf{3.} $10^7 + \num{6e7} = \num{7e7}$ \omterm{def:b1:cell-unit-of-life:cell}{células} por mL.
\textbf{4.} $7 = 2^n$: $n = 2.8$ duplicações, \qty{56}{min}.
\textbf{5.} Desligado: a lactose (a alolactose) soltou o \omterm{def:b1:expression-control:operon}{repressor}, mas
a glicose mantém o AMP cíclico baixo, de modo que a CAP não está ligada
e o promotor fica quase silencioso; além disso a permease é escassa, e
por isso entra pouca lactose.
\textbf{6.} O AMP cíclico sobe e se liga à CAP, que se liga ao promotor.
\textbf{7.} O controle negativo foi suspenso, mas faltava o controle
positivo: um promotor fraco sem CAP raramente inicia (é a \omterm{def:b1:expression-control:cap}{repressão
catabólica}).
\textbf{8.} $5000/1800 = 2.8$ tetrâmeros por segundo; $2.8\times 4096 =
\num{11400}$ resíduos por segundo.
\textbf{9.} Capacidade $20\,000\times 20 = \num{4e5}$ resíduos por
segundo: \qty{3}{\%} dos \omterm{def:b1:gene-expression:ribosome}{ribossomos}.
\textbf{10.} $5000\times 4096\times 110\times 1.66\times 10^{-24} =
\qty{3.7e-15}{g}$: \qty{2.5}{\%} da \omterm{def:b1:proteins:peptide}{proteína} da \omterm{def:b1:cell-unit-of-life:cell}{célula}.
\textbf{11.} $5000/5 = 1000$.
\textbf{12.} A lactose não atravessa a membrana sem a permease; um
mutante \emph{lacY}$^-$ tem a \omterm{def:b1:enzymes:enzyme}{enzima} mas não tem \omterm{def:b1:enzymes:enzyme}{substrato} dentro, e não
cresce em lactose (nem se induz bem, já que o indutor é feito a partir
da lactose dentro da \omterm{def:b1:cell-unit-of-life:cell}{célula}).
\textbf{13.} $1.0\times 10^{-3}\times 342 = \qty{0.342}{g/L}$, que valem
$2\times 180 = \qty{360}{g}$ de glicose por mol, ou seja,
\qty{0.36}{g/L} $= \qty{3.6e-4}{g/mL}$.
\textbf{14.} $3.6\times 10^{-4}/1.5\times 10^{-12} = \num{2.4e8}$
divisões; total final $\num{7e7} + \num{2.4e8} = \num{3.1e8}$ por mL.
\textbf{15.} $3.1/0.7 = 4.4 = 2^n$: $n = 2.1$ duplicações,
\qty{64}{min}.
\textbf{16.} Glicose \qty{56}{min}, latência \qty{30}{min}, lactose
\qty{64}{min}: \qty{150}{min} ao todo.
\textbf{17.} A lactose precisa antes ser hidrolisada, uma etapa
enzimática a mais cuja velocidade limita a oferta de glicose; e a
permease usa o gradiente de prótons para importá-la, um custo que o
transporte de glicose não paga.
\textbf{18.} $I^-$: \omterm{def:b1:enzymes:enzyme}{enzima} em \num{5000} o tempo todo (constitutiva,
depois que a glicose acaba; em glicose a CAP ainda a limita, de modo que
o nível é intermediário); a latência dele é mais curta ou inexistente,
de modo que a curva mostra uma pausa menor; o custo o retarda um pouco
em glicose.
\textbf{19.} CAP$^-$: cresce em glicose até $\num{7e7}$ e então para: o
\omterm{def:b1:expression-control:operon}{óperon lac} não pode ser ativado, a \omterm{def:b1:enzymes:enzyme}{enzima} fica perto de 5 moléculas, e a
lactose nunca é usada.
\textbf{20.} Só em lactose: nenhuma diferença no nível de \omterm{def:b1:enzymes:enzyme}{enzima} uma vez
induzida (os dois plenamente ligados) e praticamente nenhuma no
crescimento; ao $O^c$ apenas falta o atraso da indução. Só em glicose: o
$O^c$ fabrica a \omterm{def:b1:enzymes:enzyme}{enzima} inutilmente (até onde a CAP permite), e o tipo
selvagem não.
\textbf{21.} $5000\times 4096\times 4 = \num{8.2e7}$ \omterm{def:b1:water-small-molecules:atp}{ATP}: \qty{0.8}{\%}
de $10^{10}$.
\textbf{22.} Tempo de geração \qty{0.8}{\%} maior: velocidade de
crescimento \qty{0.8}{\%} menor, $\delta = 0.008$ duplicação por
geração.
\textbf{23.} $r = 2^{-0.008} = 0.99447$; $r^n = 0.111$: $n =
\ln 0.111/\ln 0.99447 = 2.20/0.00555 = 400$ gerações.
\textbf{24.} Em lactose a \omterm{def:b1:enzymes:enzyme}{enzima} é necessária de qualquer modo, de modo
que o mutante nada paga e ainda ganha a latência: ele não é selecionado
contra, e às vezes é selecionado a favor; e as linhagens de laboratório
são cultivadas por conveniência, e não em competição.
\textbf{25.} Fator de indução de 1000 (de 5 a \num{5000} moléculas); a
síntese desnecessária custa \qty{0.8}{\%} do orçamento de \omterm{def:b1:water-small-molecules:atp}{ATP}; o
experimento dura cerca de duas horas e meia.
\end{solution}
